\documentclass[aspectratio=169]{beamer}

% Shared Beamer settings for Elements of AIML lectures (UPES).
% Keep this file free of \documentclass to allow reuse via \input.

% Allow lecture sources to use shared theme/assets without copying them into each lecture folder.
% NOTE: These paths are relative to `Unit-xx/Lecture-yy/latex/` (and `Lab/Experiment-xx/latex/`).
\makeatletter
\def\input@path{{../../../_shared/}{../../../_shared/UPES-Beamer-Theme/}}
\makeatother

\usetheme{UPES}
\setbeamertemplate{navigation symbols}{}

\usepackage{amsmath}
\usepackage{amssymb}
\usepackage{booktabs}
\usepackage{graphicx}
\usepackage{hyperref}
\usepackage{xcolor}
\usepackage{listings}

% Code listings (general-purpose).
\lstset{
  basicstyle=\ttfamily\small,
  keywordstyle=\color{blue},
  commentstyle=\color{gray},
  stringstyle=\color{teal},
  showstringspaces=false,
  columns=fullflexible,
  breaklines=true
}

\usepackage{tikz}
\usetikzlibrary{arrows.meta,positioning,shapes.geometric,calc}

% Per-slide source line (references shown on the slide itself).
\newcommand{\slidesources}[1]{%
  \vfill
  {\tiny\color{gray}\raggedright\sloppy\parbox{\linewidth}{Sources: #1}\par}%
}

% Unit-wise source strings for this subject (used by slides/notes).
% Path is relative to the lecture `latex/` folder (the compilation CWD).
% Source strings for Elements of AIML (used by slides + notes).

\newcommand{\AIMLRepoURL}{https://github.com/tali7c/Elements-of-AIML}

\newcommand{\AIMLLabSources}{%
Provost and Fawcett, \textit{Data Science for Business}.;%
McKinney, \textit{Python for Data Analysis}.;%
WEKA Documentation (Waikato Environment for Knowledge Analysis), accessed 2026-02-28.;%
Matplotlib and Seaborn documentation, accessed 2026-02-28.%
}

\newcommand{\AIMLUnitISources}{%
Rich and Knight, \textit{Artificial Intelligence}, McGraw Hill, 2017.;%
Russell and Norvig, \textit{Artificial Intelligence: A Modern Approach} (4e), Ch. 1--2 (Introduction, Intelligent Agents).;%
Poole and Mackworth, \textit{Artificial Intelligence: Foundations of Computational Agents} (2e), selected sections.%
}

\newcommand{\AIMLUnitIISources}{%
Russell and Norvig, \textit{Artificial Intelligence: A Modern Approach} (4e), Ch. 7--9 (Logical Agents, First-Order Logic, Inference).;%
Huth and Ryan, \textit{Logic in Computer Science} (2e), selected sections on propositional and first-order logic.;%
Genesereth and Nilsson, \textit{Logical Foundations of Artificial Intelligence}, selected chapters.%
}

\newcommand{\AIMLUnitIIISources}{%
Mueller and Massaron, \textit{Machine Learning for Dummies}, For Dummies, 2016.;%
Mitchell, \textit{Machine Learning}, selected chapters on learning basics and evaluation.;%
Scikit-learn documentation, accessed 2026-02-28.%
}

\newcommand{\AIMLUnitIVSources}{%
Mueller and Massaron, \textit{Machine Learning for Dummies}, For Dummies, 2016.;%
James, Witten, Hastie, and Tibshirani, \textit{An Introduction to Statistical Learning} (2e), selected chapters.;%
Scikit-learn documentation, accessed 2026-02-28.%
}

\newcommand{\AIMLUnitVSources}{%
NIST, \textit{AI Risk Management Framework (AI RMF 1.0)}, 2023.;%
Barocas, Hardt, and Narayanan, \textit{Fairness and Machine Learning} (online book), selected sections.;%
Knaflic, \textit{Storytelling with Data}, selected chapters on communicating results.%
}



\title[Elements of AIML]{Elements of AIML}
\subtitle{Unit 05 -- Lecture 02: ML Applications in Banking, Security \& Healthcare}
\author{Tofik Ali}
\institute{School of Computer Science, UPES Dehradun}
\date{\today}

\graphicspath{{../images/}{../../../_shared/}{../../../_shared/UPES-Beamer-Theme/}}

\begin{document}

\UPESCoverFrame
\setcounter{framenumber}{0}

\begin{frame}
  \titlepage
  \vspace{0.3em}
  \begin{center}
    \footnotesize Topic focus: application case studies + risks + evaluation\\[0.25em]
    \footnotesize Repository: \texttt{\AIMLRepoURL}
  \end{center}
  \slidesources{\AIMLUnitVSources}
\end{frame}

\begin{frame}{Quick Links}
  \centering
  \hyperlink{sec:bank}{\beamerbutton{Banking}}\hspace{0.6em}
  \hyperlink{sec:sec}{\beamerbutton{Security}}\hspace{0.6em}
  \hyperlink{sec:health}{\beamerbutton{Healthcare}}\hspace{0.6em}
  \hyperlink{sec:risks}{\beamerbutton{Risks}}\hspace{0.6em}
  \hyperlink{sec:summary}{\beamerbutton{Summary}}
\end{frame}

\begin{frame}{Agenda}
  \tableofcontents
\end{frame}

\section{Banking}
\label{sec:bank}

\begin{frame}{Learning Outcomes}
  By the end of this lecture, you should be able to:
  \begin{itemize}[<+->]
    \item List typical ML use-cases in banking, security, and healthcare.
    \item Identify common data sources and labels in these domains.
    \item Explain domain-specific risks and why evaluation choices matter.
    \item Describe concept drift and why monitoring is needed.
  \end{itemize}
\end{frame}

\begin{frame}{Abbreviations \& Keywords}
  \begin{itemize}[<+->]
    \item AML: Anti-Money Laundering \hfill KYC: Know Your Customer
    \item PII: personally identifiable information
    \item EHR: electronic health record \hfill PHI: protected health information
    \item Drift: changing patterns over time \hfill Monitoring: ongoing checks
  \end{itemize}
\end{frame}

\begin{frame}{Case-Study Lens for Every Domain}
  \begin{itemize}[<+->]
    \item What is the \textbf{task}? (classification, regression, anomaly detection)
    \item What \textbf{data} and \textbf{labels} are available?
    \item Which \textbf{metric} matches the real error cost?
    \item What is the main \textbf{risk}? (drift, bias, adversary, safety)
    \item Which \textbf{monitoring signal} would reveal trouble early?
  \end{itemize}
\end{frame}

\begin{frame}{Banking: Typical ML Tasks}
  \begin{itemize}[<+->]
    \item Fraud detection (rare-event classification)
    \item Credit scoring (risk prediction)
    \item Customer churn prediction
    \item Customer support (NLP)
  \end{itemize}
  \vspace{0.4em}
  Key issues: imbalance, costs of false positives/negatives, fairness and regulation.
\end{frame}

\begin{frame}{Banking: Data and Labels (Examples)}
  \begin{itemize}[<+->]
    \item Transaction logs: amount, merchant, time, location, device
    \item Customer profile: age, income band, account history
    \item Labels: ``fraud'' confirmed, ``default'' event, churn yes/no
  \end{itemize}
  \small Labels may arrive late (fraud confirmed days later), affecting evaluation and monitoring.
\end{frame}

\section{Security}
\label{sec:sec}

\begin{frame}{Security: Typical ML Tasks}
  \begin{itemize}[<+->]
    \item Spam/phishing detection
    \item Intrusion detection (anomaly detection)
    \item Malware classification
  \end{itemize}
  \vspace{0.4em}
  Key issues: adversarial behavior, evolving threats (drift), false alarm fatigue.
\end{frame}

\begin{frame}{Security: What makes it different?}
  \begin{itemize}[<+->]
    \item The data distribution changes because attackers adapt (adversarial drift).
    \item FP cost: alert fatigue; FN cost: breach goes undetected.
    \item Evaluation often needs time-based splits and realistic ``red-team'' testing.
  \end{itemize}
\end{frame}

\section{Healthcare}
\label{sec:health}

\begin{frame}{Healthcare: Typical ML Tasks}
  \begin{itemize}[<+->]
    \item Risk prediction (readmission, complications)
    \item Imaging assistance (classification/segmentation)
    \item Triage and decision support
  \end{itemize}
  \vspace{0.4em}
  Key issues: safety, high cost of false negatives, data privacy, validation and clinical workflow.
\end{frame}

\begin{frame}{Healthcare: Data and Workflow Constraints}
  \begin{itemize}[<+->]
    \item Data: EHR, lab results, imaging, vitals (often missing/noisy)
    \item Strong privacy and consent constraints (PHI)
    \item Clinical validation: model must fit workflow and be safe to use
  \end{itemize}
  \small ``High accuracy'' is not enough if the model is unsafe or unusable.
\end{frame}

\section{Risks and monitoring}
\label{sec:risks}

\begin{frame}{Cross-Domain Risks (Common Themes)}
  \begin{itemize}[<+->]
    \item \textbf{Imbalance}: fraud and rare diseases are rare events.
    \item \textbf{Drift}: patterns change (new fraud strategies, new malware).
    \item \textbf{Bias/fairness}: historical decisions can create unfair models.
    \item \textbf{Privacy}: sensitive data; strict governance required.
  \end{itemize}
\end{frame}

\begin{frame}{Concept Drift (Why Monitoring Matters)}
  \begin{itemize}[<+->]
    \item Data distribution changes over time.
    \item Model performance degrades silently.
    \item Monitoring: track data statistics + performance metrics + alerting.
  \end{itemize}
\end{frame}

\begin{frame}{Monitoring Signals (Examples)}
  \begin{itemize}[<+->]
    \item Feature distribution shift (mean/std, missing rates)
    \item Model score distribution (sudden changes can indicate drift)
    \item Business metrics: chargebacks, false declines, incident volume
    \item Periodic ground-truth performance once labels arrive
  \end{itemize}
\end{frame}

\begin{frame}{Cross-Domain Comparison}
  \scriptsize
  \begin{tabular}{p{1.9cm}p{2.3cm}p{2.2cm}p{4.8cm}}
    \toprule
    \textbf{Domain} & \textbf{Task} & \textbf{Metric} & \textbf{Main risk / monitoring} \\
    \midrule
    Banking & fraud / credit risk & precision-recall + cost & proxy bias, delayed labels, drift / chargebacks, false declines \\
    Security & phishing / intrusion & recall + alert quality & adversarial drift / alert volume, confirmed incidents \\
    Healthcare & risk / imaging / triage & recall + calibration & safety under dataset shift / site-wise performance, override rates \\
    \bottomrule
  \end{tabular}
\end{frame}

\section{Summary}
\label{sec:summary}

\begin{frame}{Summary}
  \begin{itemize}
    \item Applications differ, but share patterns: classification, regression, anomaly detection.
    \item Domain risks: imbalance, drift, adversaries, safety and privacy constraints.
    \item Reliable deployment requires monitoring and careful metric choices.
  \end{itemize}
  \slidesources{\AIMLUnitVSources}
\end{frame}

\begin{frame}{Exit Questions}
  \begin{enumerate}
    \item Why is class imbalance common in fraud detection?
    \item Give one example of concept drift in security.
    \item In healthcare, why can false negatives be more costly than false positives?
    \item List two monitoring signals you would track after deployment.
  \end{enumerate}
\end{frame}

\end{document}
